%!TEX root = ../../note.tex
% Source: ../note-course/chapters/u3/s1.tex
\subsection{Convergent sequences}
\begin{defi}[Sequence]
  A \term{sequence} of real numbers is a function $X : \N \to \R$. We write
  $x_n = X(n)$ for its $n$th term and $\{x_n\}_{n=1}^{\infty}$, or simply
  $\{x_n\}$, for the sequence.
\end{defi}

\begin{defi}[Convergence]
  A sequence $\{x_n\}$ \term{converges} to $a \in \R$ if for every
  $\epsilon > 0$ there is $n_\epsilon \in \N$ such that
  \[
    n \geq n_\epsilon \implies \abs{x_n - a} < \epsilon.
  \]
  We then write $\lim_{n \to \infty} x_n = a$, or $x_n \to a$, and call $a$
  the \term{limit} of the sequence.
\end{defi}
In symbols,
\[
  x_n \to a \iff \forall \epsilon > 0\ \ \exists n_\epsilon \in \N\ \ \forall n \geq n_\epsilon,\ \ \abs{x_n - a} < \epsilon.
\]
Intuitively, for every tolerance $\epsilon > 0$, all sufficiently late terms
lie in $(a - \epsilon, a + \epsilon)$. The index $n_\epsilon$ may depend on
$\epsilon$, and once $n \geq n_\epsilon$ the inequality must hold for
\emph{every} later term, even if earlier terms lie outside the interval. For
$x_n = 1/n$ and $\epsilon = 0.3$, every term from $n_\epsilon = 4$ on lies in
the shaded band around $a = 0$:
\begin{center}
  \begin{tikzpicture}[x=0.55cm, y=2.4cm, >=Stealth]
    \fill[red!8] (0, 0) rectangle (11, 0.3);
    \draw[densely dashed, red!75!black] (0, 0.3) -- (11, 0.3) node [right] {$a + \epsilon$};
    \draw[->] (0, 0) -- (11.5, 0) node [right] {$n$};
    \draw[->] (0, 0) -- (0, 1.15) node [above] {$x_n$};
    \node [left] at (0, 0) {$a = 0$};
    \draw[densely dashed, blue!70!black] (4, 0) -- (4, 0.95);
    \node [below] at (4, 0) {$n_\epsilon = 4$};
    \foreach \n in {1, ..., 10} {
      \fill[teal!70!black] (\n, {1/\n}) circle (1.7pt);
    }
    \node [teal!70!black, anchor=west] at (1.3, 0.85) {$x_n = 1/n$};
  \end{tikzpicture}
\end{center}

\begin{defi}[$\epsilon$-neighbourhood]
  For $a \in \R$ and $\epsilon > 0$, the \term{$\epsilon$-neighbourhood} of
  $a$ is
  \[
    V_\epsilon(a) = \{x \in \R : \abs{x - a} < \epsilon\} = (a - \epsilon, a + \epsilon).
  \]
\end{defi}
So $x_n \to a$ precisely when for every $\epsilon > 0$ there is $n_\epsilon$
with $x_n \in V_\epsilon(a)$ for all $n \geq n_\epsilon$.

\begin{eg}
  $\lim_{n \to \infty} \frac{1}{n} = 0$. Given $\epsilon > 0$, the
  Archimedean property gives $n_\epsilon \in \N$ with
  $1/n_\epsilon < \epsilon$, and then for every $n \geq n_\epsilon$,
  \[
    \abs{\frac{1}{n} - 0} = \frac{1}{n} \leq \frac{1}{n_\epsilon} < \epsilon.
  \]
\end{eg}

\begin{eg}
  $\lim_{n \to \infty} \frac{4n + 3}{n + 1} = 4$. The idea is to simplify
  the distance first:
  $\abs{\frac{4n + 3}{n + 1} - 4} = \frac{1}{n + 1} < \frac{1}{n}$, so it is
  enough that $\frac{1}{n} < \epsilon$. Given $\epsilon > 0$, choose
  $n_\epsilon \in \N$ with $1/n_\epsilon < \epsilon$; then for every
  $n \geq n_\epsilon$,
  \[
    \abs{\frac{4n + 3}{n + 1} - 4} = \frac{1}{n + 1} \leq \frac{1}{n} \leq \frac{1}{n_\epsilon} < \epsilon.
  \]
\end{eg}

\subsubsection*{Uniqueness of the limit}
\begin{nprop}\label{prop:arbitrarily-close}
  If $x, y \in \R$ and $\abs{x - y} < \epsilon$ for every $\epsilon > 0$,
  then $x = y$.
\end{nprop}

\begin{proof}
  If $x \neq y$, then $\epsilon = \abs{x - y}/2 > 0$ but
  $\abs{x - y} > \epsilon$.
\end{proof}

\begin{nthm}\label{thm:limit-unique}
  A sequence in $\R$ has at most one limit.
\end{nthm}
The idea is to pass from $a$ to $b$ through a late term $x_n$, making each
of the two distances less than $\epsilon/2$.

\begin{proof}
  Suppose $x_n \to a$ and $x_n \to b$, and let $\epsilon > 0$. Choose
  $n_1, n_2 \in \N$ with $\abs{x_n - a} < \epsilon/2$ for $n \geq n_1$ and
  $\abs{x_n - b} < \epsilon/2$ for $n \geq n_2$. For
  $n \geq \max\{n_1, n_2\}$, the triangle inequality gives
  \[
    \abs{a - b} \leq \abs{a - x_n} + \abs{x_n - b} < \frac{\epsilon}{2} + \frac{\epsilon}{2} = \epsilon.
  \]
  As $\epsilon > 0$ was arbitrary, $a = b$ by
  Proposition~\ref{prop:arbitrarily-close}.
\end{proof}

\subsubsection*{Divergence and boundedness}
\begin{defi}
  A sequence \term{diverges} if it does not converge to any real number.
\end{defi}

\begin{eg}
  The sequence $0, 1, 0, 1, \ldots$ ($x_n = 0$ for odd $n$, $x_n = 1$ for
  even $n$) diverges: both $0$ and $1$ keep appearing arbitrarily late, so the
  terms cannot all stay close to one limit. Indeed, suppose $x_n \to a$ and
  take $\epsilon = 1/2$. Choose an odd and an even index beyond $n_\epsilon$.
  The odd term gives $\abs{a} < 1/2$, i.e.\ $-1/2 < a < 1/2$, and the even
  term gives $\abs{1 - a} < 1/2$, i.e.\ $1/2 < a < 3/2$. These intervals do
  not overlap.
\end{eg}

\begin{defi}
  A sequence $\{x_n\}$ is \term{bounded} if there is $M \in \R$ with
  $\abs{x_n} \leq M$ for all $n \in \N$.
\end{defi}

\begin{nthm}\label{thm:convergent-bounded}
  Every convergent sequence of real numbers is bounded.
\end{nthm}
The idea is that the tail stays within $1$ of the limit, and only finitely
many terms come before it.

\begin{proof}
  Suppose $x_n \to a$. With $\epsilon = 1$, choose $n_1$ such that
  $\abs{x_n - a} < 1$ for all $n \geq n_1$. For these $n$,
  $\abs{x_n} \leq \abs{x_n - a} + \abs{a} < 1 + \abs{a}$. The remaining
  terms form a finite set, so
  \[
    M = \max\{\abs{x_1}, \ldots, \abs{x_{n_1 - 1}}, 1 + \abs{a}\}
  \]
  bounds every term (if $n_1 = 1$, take $M = 1 + \abs{a}$).
\end{proof}
The converse fails: $0, 1, 0, 1, \ldots$ is bounded but diverges.

\begin{eg}
  The following sequences are unbounded: for every $M > 0$ some term has
  absolute value greater than $M$.
  \begin{enumerate}
    \item $a_n = \frac{n^2 + 1}{n}$: here $\abs{a_n} = n + \frac{1}{n}$, and
      choosing $n_M \in \N$ with $n_M > M$ gives $\abs{a_{n_M}} > M$.
    \item $a_n = (-1)^n n$: here $\abs{a_n} = \abs{(-1)^n}\abs{n} = n$, which
      exceeds $M$ as soon as $n > M$.
  \end{enumerate}
  \begin{own}
    By Theorem~\ref{thm:convergent-bounded}, neither sequence converges.
  \end{own}
\end{eg}

\subsubsection*{Divergence to \texorpdfstring{$\pm\infty$}{plus or minus infinity}}
\begin{defi}
  A sequence $\{x_n\}$ \term{diverges to $+\infty$}, written
  $\lim_{n \to \infty} x_n = +\infty$, if for every $M > 0$ there is
  $n_M \in \N$ such that $x_n > M$ for all $n \geq n_M$. It
  \emph{diverges to $-\infty$}, written $\lim_{n \to \infty} x_n = -\infty$,
  if for every $M > 0$ there is $n_M \in \N$ such that $x_n < -M$ for all
  $n \geq n_M$.
\end{defi}

\begin{warning}
  Although we write $\lim_{n \to \infty} x_n = \pm\infty$, such a sequence
  is still divergent: it does not converge to any real number.
\end{warning}

\begin{eg}\leavevmode
  \begin{enumerate}
    \item $\lim_{n \to \infty} \frac{n^2}{1 - n} = -\infty$.
    \item $\lim_{n \to \infty} \bigl(\sqrt{n} + 1\bigr) = +\infty$.
  \end{enumerate}
  The proofs were not completed in lecture. For (ii), the suggested approach
  was to choose $n$ so large that $\sqrt{n} > M - 1$.
  % TODO(check): complete both proofs once the next lecture covers them.
\end{eg}

\begin{thm}
  Suppose $\lim_{n \to \infty} x_n = a$ and $\lim_{n \to \infty} y_n = b$.
  Then for every $c \in \R$:
  \begin{enumerate}
    \item the constant sequence $c, c, c, \ldots$ converges to $c$;
    \item $\lim_{n \to \infty} c x_n = ca$.
  \end{enumerate}
\end{thm}
% TODO(check): the lecture stated this without proof so far; the hypothesis
% on y_n is not used by (i) or (ii) -- presumably later parts involve it.
